\documentclass[11pt]{article}

\usepackage[margin=1in]{geometry}
\usepackage{amsmath,amssymb,amsthm}
\usepackage{booktabs}
\usepackage{hyperref}

\hypersetup{
  colorlinks=true,
  linkcolor=blue,
  citecolor=blue,
  urlcolor=blue
}

\newtheorem{theorem}{Theorem}
\newtheorem{proposition}{Proposition}
\newtheorem{definition}{Definition}
\newtheorem{remark}{Remark}

\title{NNR9 Draft: Lean4-Audited Numerical Readout for the Standing Neutrino Splitting Ratio}
\author{Draft prepared from the AASC/Lean4 development}
\date{\today}

\begin{document}

\maketitle

\begin{abstract}
This note records the present AASC/Lean4 status of the numerical readout
extension of the standing neutrino splitting-ratio program.  The earlier
endpoint identified the admissible same-scope branch as a singleton but did
not evaluate that singleton numerically.  The new step builds an AASC
standing-admission certificate for the same-scope source-row interface of the
standing neutrino target.  External material is used only as cited provenance
and row data; it is not treated as the structural proof authority.  From the
AASC standing-admission certificate, the Lean4 development constructs the
same-scope row admission object, then constructs a bundled readout certificate
and proves rational decimal windows for the AASC standing coordinate, for the
inverse hierarchy coordinate printed by the source, and for the shifted
hierarchy coordinate.

The Lean4 pipeline is formalized in
\texttt{MaleyLean.Papers.Neutrino.SourceTheorems.MathlibSpectralOperator}.
In the local workspace, \texttt{lake build MaleyLean.Papers.Neutrino} passes.
The file scan found no \texttt{sorry}, \texttt{admit}, \texttt{axiom}, or
\texttt{unsafe} occurrence in the relevant Lean file.  The broad row-admission
socket is now discharged from the stronger corpus-shaped certificate
\texttt{AASCNeutrinoSourceRowStandingAdmission}; the remaining audit burden is
therefore explicit in that certificate's fields rather than hidden inside a
single unnamed premise.
\end{abstract}

\section{Purpose}

The previous NNR endpoint may be summarized as follows:
\[
  B_{\rho}^{\mathrm{adm}} = \{\rho_{\nu}^{\mathrm{current}}\}.
\]
This is a branch-singleton result.  It says that the current standing ratio is
the unique admissible same-scope branch, but it does not by itself supply a
decimal value.  The remaining problem was the construction of a lawful
numerical evaluator for the unique surviving branch.

The present draft records the AASC standing-admission evaluator route.  The
evaluator is not built by selecting an empirical best-fit value and not by
defining the root to be the gap quotient.  Instead, AASC supplies the
same-scope standing certificate; the cited external package supplies
provenance and row data; Lean transports the admitted row law through the
existing AASC readout machinery.

\section{External Source Provenance}

The external source package used for provenance and row data is:
\begin{quote}
Mehrdad Pajuhaan, \emph{A closed, fit-free prediction for the neutrino mass
hierarchy from the charged-lepton emergent pipeline}, Zenodo,
DOI: \href{https://doi.org/10.5281/zenodo.17169926}{10.5281/zenodo.17169926}.
\end{quote}

The package contains:
\begin{itemize}
  \item \texttt{Mass-Hierarchy-Neutrinos.pdf}
  \item \texttt{Neutrinos Mass Hierarchy.py}
\end{itemize}

The companion Python source states that the coefficient pipeline does not use
neutrino masses or mass-squared splittings to set the coefficients.  In the
local audit, the relevant lines include the contamination audit, the
\(P_{1/R}\) definition, the \(k=1,3,5\) row construction, and the printed
hierarchy readouts.

\section{Coordinate Convention}

The AASC standing coordinate is
\[
  \rho = \frac{\Delta m^2_{21}}{\Delta m^2_{31}},
\]
that is, the solar-over-atmospheric splitting ratio.

The external source prints the inverse hierarchy coordinate
\[
  R_{31} = \frac{\Delta m^2_{31}}{\Delta m^2_{21}},
\]
and also the shifted coordinate
\[
  R_{3\ell} = R_{31} - \frac{1}{2}.
\]
Thus the printed hierarchy coordinate is not the same coordinate as the AASC
standing ratio.  The bridge is
\[
  \rho = R_{31}^{-1}.
\]

This distinction is formalized in Lean by:
\[
\begin{array}{ll}
\texttt{aascNeutrinoR31\_eq\_R3ell\_plus\_half}, &
\texttt{aascNeutrinoStandingRatioReadout\_eq\_inv\_R31},\\
\texttt{aascNeutrinoStandingRatioReadout\_eq\_inv\_R3ell\_plus\_half}.&
\end{array}
\]

\section{Source Row Law}

The AASC source-row law used in the Lean translation is:
\[
  P_{1/R}(k) = 1 + \beta k^2 + \delta k^4,
  \qquad k \in \{1,3,5\}.
\]
The finite-decimal source readback is:
\[
  \beta = 0.0000780161359733586,
  \qquad
  \delta = 0.0119354363008418,
  \qquad
  s = 1.
\]

The mass-squared row expression used by Lean is:
\[
  L(k) = k^4 P_{1/R}(k).
\]
The corresponding Lean objects are:
\[
\begin{array}{ll}
\texttt{aascNeutrinoBeta},&
\texttt{aascNeutrinoDelta},\\
\texttt{aascNeutrinoScale},&
\texttt{aascNeutrinoP1OverR},\\
\texttt{aascNeutrinoMassSquaredLevel},&
\texttt{aascNeutrinoNumeralParameterTable}.
\end{array}
\]

\section{Lean4 Formalization Status}

The formalization is in:
\[
  \texttt{MaleyLean.Papers.Neutrino.SourceTheorems.MathlibSpectralOperator}.
\]

The local build command is:
\[
  \texttt{lake build MaleyLean.Papers.Neutrino}.
\]
This build succeeds in the current workspace.  The relevant Lean file was
searched for \texttt{sorry}, \texttt{admit}, \texttt{axiom}, and
\texttt{unsafe}; no such proof placeholders or extra axiom declarations were
found.  The only reported Lean output is a pre-existing unrelated linter
warning recommending \texttt{simp} instead of \texttt{simpa} at one location.

\begin{remark}
This statement concerns the local Lean4 file and build route.  It does not
assert that the external physics source is itself proof authority.  The
external source provides row data and provenance; the AASC admission of those
rows is represented by an explicit Lean certificate whose fields must be
audited against the AASC corpus.
\end{remark}

\section{AASC Admission Interface}

The broad row-admission object is:
\[
  \texttt{AASCNeutrinoReadoutSameScopeAdmissionRows}.
\]
This object states that the AASC standing three-flavor algebra reads back to
the AASC-admitted source rows at \(k=1,3,5\), and that the existing QNF,
singleton, and ordering/audit guards are active.

In the present version this object is produced from the stronger AASC
standing-admission certificate:
\[
  \texttt{AASCNeutrinoSourceRowStandingAdmission}.
\]
That certificate names the finite PMNS invariant-witness scope,
mass-squared-splitting witness standing, same-scope rival exhaustion,
bridge-anchor admissibility, evidence-ledger completeness, same-standing
carrier readback, quotient/skin transport compatibility, and source-row
non-empiricality conditions.  The Lean discharge theorem is:
\[
  \texttt{AASCNeutrinoSourceRowStandingAdmission.discharge\_sameScopeAdmissionRows}.
\]

From the produced admission object, Lean constructs:
\[
\begin{array}{ll}
\texttt{aascNeutrinoSameScopeAdmissionRowsAsPrimitiveEndpointRows},\\
\texttt{aascNeutrinoSameScopeAdmissionRowsAsSingletonEndpoint},\\
\texttt{aascNeutrinoSameScopeAdmissionRows\_hasCurrentStandingEvaluation}.
\end{array}
\]

The final bundled AASC certificate is:
\[
  \texttt{AASCNeutrinoNumericalReadoutCertificate}.
\]

\section{Numerical Readout}

The source-row computation gives the following certified rational windows:
\[
  0.0299486814926010
  \leq \rho
  \leq
  0.0299486814926012,
\]
\[
  33.3904516045907
  \leq R_{31}
  \leq
  33.3904516045908,
\]
\[
  32.8904516045907
  \leq R_{3\ell}
  \leq
  32.8904516045908.
\]

In Lean these are exposed by:
\[
\begin{array}{ll}
\texttt{aascNeutrinoSourceGapRatio\_in\_decimalWindow},\\
\texttt{aascNeutrinoSourceInverseHierarchyRatio\_in\_decimalWindow},\\
\texttt{aascNeutrinoSourceR3ellRatio\_in\_decimalWindow}.
\end{array}
\]

After same-scope AASC admission, the selected root inherits these windows:
\[
\begin{array}{ll}
\texttt{aascNeutrinoNumericalReadoutCertificate\_selectedRoot\_in\_decimalWindow},\\
\texttt{aascNeutrinoNumericalReadoutCertificate\_selectedRoot\_inverse\_in\_decimalWindow},\\
\texttt{aascNeutrinoNumericalReadoutCertificate\_selectedRoot\_r3ell\_in\_decimalWindow}.
\end{array}
\]

\section{Main Endpoint}

\begin{theorem}[AASC standing-admission numerical readout]
Assume the AASC standing source-row certificate
\[
  \texttt{AASCNeutrinoSourceRowStandingAdmission}.
\]
This certificate packages the finite PMNS invariant-witness standing,
mass-squared splitting witness standing, same-scope rival exhaustion,
bridge-anchor admissibility, evidence-ledger completeness, same-standing
source-row carrier, quotient/skin transport compatibility, and the
non-empirical source-row provenance.  From it, Lean constructs
\[
  \texttt{AASCNeutrinoReadoutSameScopeAdmissionRows}
\]
and then proves:
\[
  0.0299486814926010
  \leq \rho_{\nu}^{\mathrm{current}}
  \leq
  0.0299486814926012,
\]
\[
  33.3904516045907
  \leq
  \left(\rho_{\nu}^{\mathrm{current}}\right)^{-1}
  \leq
  33.3904516045908,
\]
and
\[
  32.8904516045907
  \leq
  \left(\rho_{\nu}^{\mathrm{current}}\right)^{-1} - \frac{1}{2}
  \leq
  32.8904516045908.
\]
\end{theorem}

The corresponding Lean theorem is:
\[
  \texttt{AASCNeutrinoSourceRowStandingAdmission.finalNumericalReadout}.
\]

\section{Theorem Map}

\begin{center}
\begin{tabular}{p{0.44\linewidth}p{0.48\linewidth}}
\toprule
Manuscript claim & Lean object \\
\midrule
Source provenance & \texttt{AASCNeutrinoCoefficientPipelineSourceDerived} \\
No empirical coefficient import & \texttt{AASCNeutrinoNoMassOrSplittingInputInCoefficientPipeline} \\
No hidden selected-root encoding & \texttt{AASCNeutrinoNoHiddenSelectedRootEncodingInSourceRows} \\
Concrete coefficient table & \texttt{aascNeutrinoNumeralParameterTable} \\
AASC row law & \texttt{aascNeutrinoP1OverR} and \texttt{aascNeutrinoMassSquaredLevel} \\
Coordinate inversion & \texttt{aascNeutrinoStandingRatioReadout\_eq\_inv\_R31} \\
Source \(\rho\) window & \texttt{aascNeutrinoSourceGapRatio\_in\_decimalWindow} \\
Source \(R_{31}\) window & \texttt{aascNeutrinoSourceInverseHierarchyRatio\_in\_decimalWindow} \\
Source \(R_{3\ell}\) window & \texttt{aascNeutrinoSourceR3ellRatio\_in\_decimalWindow} \\
Corpus-shaped row admission certificate & \texttt{AASCNeutrinoSourceRowStandingAdmission} \\
Discharge of broad row socket & \texttt{AASCNeutrinoSourceRowStandingAdmission.discharge\_sameScopeAdmissionRows} \\
AASC row admission object produced & \texttt{AASCNeutrinoReadoutSameScopeAdmissionRows} \\
Bundled readout certificate & \texttt{AASCNeutrinoNumericalReadoutCertificate} \\
Final numerical theorem & \texttt{AASCNeutrinoSourceRowStandingAdmission.finalNumericalReadout} \\
\bottomrule
\end{tabular}
\end{center}

\section{Interpretation}

The result is a standing-admission numerical readout theorem.  It is stronger
than the previous branch-singleton endpoint because it supplies rational
decimal windows for the unique selected root after source-row admission.  It
is still not an unconditional theorem from the empty context: the manuscript
must exhibit or cite the AASC standing-admission certificate.  The important
upgrade is that the certificate is no longer a vague row-admission label; its
auditable fields are explicitly named in Lean.

In plain language, the current status is:
\begin{quote}
Given the AASC source-row standing certificate for the standing neutrino
splitting target, the Lean4 pipeline, without additional axioms or proof
placeholders in the relevant file, constructs the same-scope admission rows
and derives the numerical windows above for the current standing ratio and
its inverse hierarchy coordinates.
\end{quote}

\section{One-Anchor Continuation}

The numerical readout above is a ratio/coordinate anchor, not an absolute-mass
anchor.  Lean now records this distinction in a follow-on one-anchor layer.
The ratio-only residual family is represented by
\[
  m_1^2=t,\qquad
  m_2^2=t+\rho D,\qquad
  m_3^2=t+D.
\]
For \(D\neq 0\), Lean proves that this family has gap ratio \(\rho\), and
that changing \(t\) leaves the ratio unchanged.  The relevant objects are:
\[
\begin{array}{ll}
\texttt{threeFlavorRatioOffsetGapFamily},\\
\texttt{threeFlavorRatioOffsetGapFamily\_gapRatio},\\
\texttt{threeFlavorRatioOffsetGapFamily\_offsetInvariant}.
\end{array}
\]

The MR6-style continuation therefore does not claim zero-anchor absolute
neutrino masses.  It states the lawful display theorem: a full projective
mass-squared shape plus one target-aligned metric anchor reconstructs an
absolute mass-squared triple while preserving the gap ratio.  The Lean objects
are:
\[
\begin{array}{ll}
\texttt{aascNeutrinoSourceProjectiveShapeLevel},\\
\texttt{aascNeutrinoSourceProjectiveMassSquaredShape},\\
\texttt{aascNeutrinoSourceProjectiveMassSquaredShape\_gapRatio\_eq\_source},\\
\texttt{ThreeFlavorProjectiveMassSquaredShape},\\
\texttt{ThreeFlavorOneAnchorMassSquaredDisplayAnchor},\\
\texttt{AASCNeutrinoOneAnchorMetricInput},\\
\texttt{aascNeutrinoOneAnchorSpectrumReconstruction},\\
\texttt{aascNeutrinoNNR10OneAnchorMassSpectrumDisplay},\\
\texttt{aascNeutrinoOneAnchorSpectrum\_gapRatio\_eq\_NNR9GapRatio},\\
\texttt{aascNeutrinoOneAnchorSpectrum\_gapRatio\_eq\_selectedRoot},\\
\texttt{aascNeutrinoOneAnchorSpectrum\_anchorLevel\_eq\_anchorValue},\\
\texttt{AASCNeutrinoOneAnchorAbsoluteMassSquaredSpectrumAvailable},\\
\texttt{aascNeutrinoOneAnchorMetricInput\_discharge\_absoluteSpectrumNeed},\\
\texttt{AASCNeutrinoAtmosphericGapMetricAnchor},\\
\texttt{aascNeutrinoAtmosphericGapScale},\\
\texttt{aascNeutrinoAtmosphericGapAnchoredMassSquared},\\
\texttt{aascNeutrinoAtmosphericGapAnchored\_atmosphericGap\_eq\_anchor},\\
\texttt{aascNeutrinoAtmosphericGapAnchored\_gapRatio\_eq\_selectedRoot},\\
\texttt{AASCNeutrinoAtmosphericGapAnchoredSpectrumAvailable},\\
\texttt{aascNeutrinoAtmosphericGapMetricAnchor\_discharge\_spectrumNeed},\\
\texttt{threeFlavorOneAnchorReconstructedMassSquared},\\
\texttt{threeFlavorOneAnchorReconstructed\_anchor\_eq},\\
\texttt{threeFlavorOneAnchorReconstructed\_gapRatio\_eq\_shape},\\
\texttt{ThreeFlavorOneAnchorNeutrinoSpectrumReconstruction},\\
\texttt{AASCNNR10OneAnchorNeutrinoMassSpectrumDisplay}.
\end{array}
\]

Thus the charged-lepton symmetry is captured in the right form: AASC can
license a one-anchor metric display once the AASC source-row projective shape
and a target-aligned anchor are admitted, while the NNR9 ratio readout alone
does not erase the offset/scale freedom.  In particular, Lean now proves that
the reconstructed one-anchor spectrum preserves the NNR9 selected root as its
gap ratio and that the anchored component equals the admitted anchor value.
For the neutrino-specific case, Lean also provides the atmospheric-gap anchor:
if an admitted \(\Delta m^2_{31}\) anchor is supplied, the scale
\[
  \lambda =
  \frac{\Delta m^2_{31}}{L_5-L_1}
\]
constructs the absolute mass-squared display \(m_i^2=\lambda L_i\), proves
that the atmospheric gap is exactly the anchor, and preserves the NNR9 selected
root as the gap ratio.

\subsection{Worked Atmospheric-Gap Display}

For a concrete display, take the normal-ordering atmospheric gap anchor
\[
  \Delta m^2_{31}=2.513\times 10^{-3}\,\mathrm{eV}^2.
\]
This is an external display anchor, cited from the NuFIT global oscillation
fit convention in which \(\Delta m^2_{3\ell}=\Delta m^2_{31}>0\) for normal
ordering; it is not an AASC zero-anchor derivation of absolute mass.

The AASC source-row projective levels are
\[
\begin{aligned}
L_1&=1.0120134524368151586,\\
L_3&=159.3652713329476282194,\\
L_5&=5288.4988071409118531250.
\end{aligned}
\]
Hence
\[
  \lambda =
  \frac{2.513\times 10^{-3}}{5288.4988071409118531250
    -1.0120134524368151586}
  \approx 4.7527305468\times 10^{-7}.
\]
The anchored display is
\[
  (m_1^2,m_2^2,m_3^2)=
  \lambda(L_1,L_3,L_5),
\]
and therefore
\[
\begin{aligned}
m_1&\approx 0.0006935\,\mathrm{eV}
      =0.694\,\mathrm{meV},\\
m_2&\approx 0.008703\,\mathrm{eV}
      =8.703\,\mathrm{meV},\\
m_3&\approx 0.050135\,\mathrm{eV}
      =50.135\,\mathrm{meV}.
\end{aligned}
\]
The corresponding sum is
\[
  \sum_i m_i \approx 0.05953\,\mathrm{eV}
  =59.53\,\mathrm{meV}.
\]
This display checks:
\[
  \Delta m^2_{31}=0.002513\,\mathrm{eV}^2,\qquad
  \Delta m^2_{21}\approx 7.526\times 10^{-5}\,\mathrm{eV}^2,
\]
and
\[
  \frac{\Delta m^2_{21}}{\Delta m^2_{31}}
  \approx 0.02994868149.
\]

\begin{thebibliography}{9}

\bibitem{ExternalZenodoSource}
M. Pajuhaan,
\emph{A closed, fit-free prediction for the neutrino mass hierarchy from the
charged-lepton emergent pipeline},
Zenodo, 2025.
DOI: \href{https://doi.org/10.5281/zenodo.17169926}{10.5281/zenodo.17169926}.
Record: \url{https://zenodo.org/records/17169926}.

\bibitem{NuFIT60}
I. Esteban, M. C. Gonzalez-Garcia, M. Maltoni, I. Martinez-Soler,
J. P. Pinheiro, and T. Schwetz,
\emph{NuFit-6.0: updated global analysis of three-flavor neutrino oscillations},
JHEP \textbf{2024}, 216 (2024).
DOI: \href{https://doi.org/10.1007/JHEP12(2024)216}
{10.1007/JHEP12(2024)216}.
NuFIT v6.0 results page:
\url{https://www.nu-fit.org/?q=node/294}.

\end{thebibliography}

\end{document}
